% Section 9.3, from lectures/L09/README.md: the objectives, the evaluation questions, and the next
% lecture.

\section{Review}
\label{c9:sec:review}

\needspace{6\baselineskip}
\subsection*{What you should be able to do now}
After this chapter, you should be able to:
\begin{itemize}
  \item \textbf{Configure the AVR SPI master registers} correctly for the transport contract, and
    explain each bit they set.
  \item \textbf{Implement \code{AvrSpi}} so that the existing driver and its host tests are reused
    unchanged.
  \item \textbf{Explain which host-side idioms do not survive a freestanding build}, and what
    replaces them.
  \item \textbf{Synthesize and program the DE0-CV}, and read a fit and timing report well enough to
    say whether the design met timing at 50\,MHz.
\end{itemize}

\needspace{6\baselineskip}
\subsection*{Questions to test yourself}
You should be able to explain:
\begin{enumerate}
  \item Which idiom from the host build fails to compile freestanding, and how does this chapter
    replace it?
  \item The AVR SPI runs at 1\,MHz: which prescaler setting produces it from a 16\,MHz clock, and
    what actually limits \code{SCK} on this bench, given that the FPGA slave could keep up with
    considerably more?
  \item If \code{SPDR} is read a cycle too early, before \code{SPIF} is set, what value comes back,
    and how does the transport avoid it?
\end{enumerate}

\needspace{6\baselineskip}
\subsection*{Looking ahead}
In \chapref{10} both halves are on the bench: two links, a level shifter, and a bring-up ladder that
ends with an echo application driven end to end.
